\section{Architecture}
\redacteur{\acompleter[nom]}

\consigne{Le schéma est obligatoire. Partez du schéma Mermaid du livrable 3, exportez-le en PDF ou PNG (par exemple sur mermaid.live) et placez-le dans \texttt{figures/}. Chaque composant du schéma est expliqué dans le texte et chaque choix technique est justifié face à au moins une alternative.}

\insererfigure{architecture.pdf}{0.95}{Architecture du système}{fig:architecture}

\subsection{Composants}

\begin{table}[H]\centering
\caption{Composants et choix techniques}\label{tab:composants}
\begin{tabularx}{\linewidth}{lllX}
\toprule
\textbf{Composant} & \textbf{Technologie} & \textbf{Alternative écartée} & \textbf{Justification} \\
\midrule
Modèle & \acompleter & \acompleter & \acompleter \\
Recherche ou mémoire & \acompleter{} & \acompleter{} & \acompleter{} (voir par exemple \citet{lewis2020rag}) \\
API & FastAPI \citep{fastapi} & \acompleter & \acompleter \\
Interface & \acompleter & \acompleter & \acompleter \\
\bottomrule
\end{tabularx}
\end{table}

\subsection{Parcours d'une requête}
\acompleter[de l'action de l'utilisateur jusqu'à la réponse, en suivant la figure~\ref{fig:architecture}]

\subsection{Agents et orchestration}
\consigne{Rôle de chaque agent, outils qu'il peut appeler, façon dont les agents échangent. Obligatoire pour une équipe de 4.}
\acompleter
